\documentclass{article}
\usepackage[T1]{fontenc}
\usepackage{lmodern}
\usepackage{graphicx}
\usepackage{amsmath,amssymb}
\usepackage[a4paper,margin=2cm]{geometry}
\usepackage[absolute,overlay]{textpos}
\setlength{\TPHorizModule}{1mm}
\setlength{\TPVertModule}{1mm}
\usepackage[indonesian]{babel}
\usepackage{hyperref}
\setlength{\emergencystretch}{2em}
% unit: Halaman 25

\begin{document}

% === Halaman 25 (python/native) ===

Kurva Eliptik pada GF(p)

• Bentuk umum kurva eliptik pada GF(p) (atau Fp) :


y2  x3 + ax + b  (mod p)

 yang dalam hal ini p adalah bilangan prima dan elemen-elemen 

himpunan di dalam medan galois adalah \{0, 1, 2, …, p – 1\}

25

\end{document}
